

\justify 

\setlength{\parindent}{0pt}


% ==================================================================================================================================
% Introduction

Dans la chapitre 1, nous avons vu comment "réduire" une application linéaire pour l'exprimer par une matrice diagonale 
sans changer son comportement intrinsèque. Ce résultat, très fort, n'est malheureusement pas banal. 
On va donc réduire la matrice représentative de notre application, non pas sous forme diagonale, mais sous forme triangulaire. 
En pratique, cela représentera un avantage moins important que la diagonalisation mais cela sera toujours préférable
comparé à des calculs avec la matrice de départ. On appelle cette opération (i.e changement de base) la trigonalisation. 

Le procédé de trigonalisation étant plus complexe que celui de diagonalisation, nous allons donc élaborer une méthode (des drapeaux)
pour pouvoir trigonaliser toute fonction respectant quelques critères. 


% ==================================================================================================================================
% Trigonalisation définitions

\section{Trigonalisation - Rappels}

Redéfinissons proprement la trigonalisation comme vu en fin de chapitre 1. 

\begin{definition}[Trigonalisation]
    Un endomorphisme $f$ représenté par une matrice A dans une base $\mathcal{B}$ est dit trigonalisable si 
    il existe une base $\mathcal{B}'$ dans laquelle il est représenté par une matrice triangulaire supérieure. 

    \vspace{0.3cm}

    De même une matrice $A$ est dite trigonalisable si elle est semblable à une matrice triangulaire supérieure. 
\end{definition}

\begin{prop}[Lien matrice/endomorphisme]
    Un endomorphime $f$ représenté par une matrice $A$ est trigonalisable ssi $A$ l'est. 
\end{prop}

\subsection{Premiers critères}

\begin{criteria}[Fondamental de Trigonalisation]
    Soient E un $\K$-espace vectoriel et $f \in \mathcal{L}(E)$. $f$ est trigonalisable ssi son polynôme caractéristique est scindé sur $\K$. 
\end{criteria}

\begin{corollary}[La puissance de $\C$]
    Tout endormorphisme est trigonalisable sur $\C$. 
\end{corollary}

\newpage 

\section{Méthode de Dunford}

\subsection{Théorie}

\begin{definition}[Sous-espace caractéristique]
    Soit $f$ un endomorphisme trigonalisable de polynôme caractéristique 
        $$ P_f = (\lambda_1 - X)^{\alpha_1} \dots (\lambda_p -X)^{\alpha_p} $$ 
    On appelle sous-espace caractéristique de $f$ associé à la valeur propre $\lambda_i$, le sous-espace vectoriel 
        $$ \boxed{ F_i = \ker((u - \lambda_i \; \text{Id})^{\alpha_i}) } $$ 
\end{definition}

\begin{remark}
    Attention à la différence entre les sous-espaces propres et les sous-espaces caractéristiques !
    Les seconds sont définis en fonction de la multiplicité de la valeur propre alors que les premiers non. 
\end{remark}

\begin{theorem}[Trigonalisation et sous-espace caractérisitique]
    Soit $f$ un endomorphisme trigonalisable
    \begin{itemize}
        \item Les $F_i$ sont stables par $f$ et :
            \[ E = \bigoplus_{i \in \{1,\dots,p\}} F_i \]
        \item Pour toute valeur propre, le sous-espace \textbf{propre} correspondant est inclus dans le sous-espace 
        \textbf{caractéristique} correspondant. 
            \[ \forall i \in \{1,\dots,p\}, \quad E_i \subset F_i \] 
        \item Pour toute valeur propre, la dimension du sous-espace \textbf{caractéristique} correspondant est égale à sa multiplicité. 
            \[ \forall i \in \{1,\dots,p\}, \quad \text{dim}(F_i) = \alpha_i \]
    \end{itemize}
\end{theorem}


\subsection{En pratique}

En résumé, la méthode pratique de trigonalisation se résume en quelques étapes. 
Soit $f$ un endomorphisme trigonalisable représenté par une matrice $A$ dans une base $\mathcal{B}'$. 
Pour trigonaliser $f$ (ou $A$), il faut :
\begin{itemize}
    \item Calculer le polynôme caractéristique $P_f = P_A$. 
    En le factorisant, en déduire le spectre de $f$ et, pour chaque valeur propre, son ordre de multiplicité. 
    Si $P_f$ est scindé, alors $f$ est trigonalisable. 
    \item Pour chaque valeur propre, on calcule le sous-espace propre associé, on en donne la dimension et une base. 
    \item Si pour chaque valeur propre on a égalité entre la dimension du sous-espace propre et la multiplicité de la valeur 
    propre, alors $f$ est diagonalisable. Sinon, on continue la trigonalisation. 
    \item Pour chaque valeur propre (sympathique) pour laquelle on a égalité entre multiplicité et dimension du sep,
    on détermine une base du sous-espace propre. Le bloc correspondant sera diagonal.
    \item Pour chaque valeur propre (pénible) pour laquelle on n’a pas égalité entre multiplicité et dimension du sep,
    on détermine le drapeau complet $E_\lambda = K_1 \subset \dots \subset K_p = F_\lambda$.
    On construit une base de $\mathcal{B}_i$ de $F_\lambda$ par la méthode des drapeaux.
    \item On calcule l'image de chaque $\mathcal{B}_i$ par $f$, ce qui fournit le bloc correspondant à la valeur propre $\lambda_i$. 
    \item Une base de trigonalisation est obtenue en faisant l'union des bases obtenues précédement. 
    \item La matrice triangulaire (de Dunford) est obtenue à partir des bloc précédents. 
\end{itemize}


\begin{example}[Trigonalisation]
    Illustrons cette méthode au travers d'un exemple. 
    Soit $A \in \mathcal{M}_3(\R)$ telle que :
        \[ A = 
            \begin{pmatrix}
                4 & 1 & 2 \\ 
                1 & 4 & 1 \\ 
                2 & 1 & 1 
            \end{pmatrix}
        \]
    Trigonalisons A avec la méthode vue prédécement. Procédons par étapes. 

    \begin{itemize}
        \item \textbf{Calcul du polynôme caractéristique et du spectre de $A$ :}
        \[
            P_A = \begin{vmatrix}
                        4-X & 1 & 2 \\ 
                        1 & 4-X & 1 \\ 
                        2 & 1 & 1-X 
                    \end{vmatrix} 
                = \dots = (X-6)(X-3)^2
        \]
        On a donc $\text{Sp}(A) = \{6,3\}$ où 3 est de multiplicité 2. 
        Le polynôme caractéristique est scindé donc $A$ est trigonalisable dans $\R$. 

        \item \textbf{Calcul des sous-espaces propres :}
            \begin{itemize}
                \item \textbf{Sous espace associé à 6 :}
                \[ A - 6 \text{Id} = 
                    \begin{pmatrix}
                        -2 & 1 & 2 \\ 
                        1 & -2 & 1 \\ 
                        2 & 1 & -2 
                    \end{pmatrix}
                \]
                Grâce au pivot de Gauss, on obtient $E_6 = \ker (A - 6 \text{Id}) = \text{vect}((1,1,1))$. 

                La dimension du sous-espace propre est égale à l'ordre de multiplicité de la valeur propre, 
                on a dont affaire à une valeur propre "sympathique". 
                \item \textbf{Sous espace associé à 3 :}
                \[ A - 3 \text{Id} = 
                    \begin{pmatrix}
                        1 & 1 & 2 \\
                        1 & 1 & 1 \\ 
                        2 & 1 & 1 
                    \end{pmatrix}
                \]
                Toujours grâce au pivot de Gauss, on obtient $E_3 = \ker (A - 3 \text{Id}) = \text{vect}(())$ 
            \end{itemize}
    \end{itemize}
\end{example}